We now address the related but quantitatively distinct problem of estimating the number of unique elements within a large-scale dataset.

\begin{definition}[The Cardinality Estimation Problem]
    Given a multiset $S$ of elements drawn from a universe $\mathcal{U}$, the cardinality estimation problem is to determine $|S|$, defined as the number of distinct elements in $S$.
\end{definition}

In a distributed, data-parallel framework, a naive implementation of this task necessitates a multi-stage execution. In an initial stage, each parallel worker processes a partition of the input multiset to produce a local set of unique elements. Subsequently, the framework must perform a shuffle operation, transmitting all unique elements from all partitions across the network to a common set of reducer nodes for final aggregation and deduplication \cite{Zaharia2012ResilientDD}.

This shuffle operation constitutes the primary bottleneck. The communication cost of this approach is not fixed; it is proportional to the number of unique elements and their size. For datasets with high cardinality, the network I/O required for the shuffle dominates the total execution time, rendering the approach impractical for interactive or large-scale analytics.

The challenge here is to enable cardinality aggregation without the shuffle of raw data. A distributed cardinality estimation oracle must support the combination of partial estimates computed on different data partitions into a global estimate, where this combination operation satisfies three requirements.

First, produces a result with error bounds identical to those of an oracle constructed in a single pass over the complete dataset. This ensures that distributed aggregation does not introduce additional approximation error beyond the inherent error of the probabilistic estimator.

Second, executes in time independent of the dataset cardinality. The combination of two partial estimates must require computational cost proportional only to the size of the estimate representations, not to the number of elements they summarize. This transforms the communication-bound shuffle operation into a fixed-cost sketch transmission.

Third, can be performed in any order without coordination between nodes. The final global estimate must be identical regardless of the sequence in which partial estimates are combined. This enables distributed reduction without sequential dependencies or central coordination bottlenecks.

We observe that these three requirements collectively specify a class of binary operations over cardinality estimation oracles. The requirement for a property that satisfies these conditions motivates the subsequent paragraphs.

\subsection{Composability}

The constraints identified in the preceding sections impose requirements on the structure of a distributed cardinality estimation oracle. We now introduce the property that resolves the shuffle bottleneck constraint.

\begin{definition}[Composability] \label{def:composability}
    Let $\mathcal{D}_1, \dots, \mathcal{D}_k$ be $k$ instances of a probabilistic oracle, where each $\mathcal{D}_i$ is constructed over a data partition $S_i$. The oracle exhibits composability if there exists a binary operation $\cup$ such that the combined oracle $\mathcal{D}_{agg} = \mathcal{D}_1 \cup \dots \cup \mathcal{D}_k$ produces estimates with expected relative error $\mathbb{E}[|\hat{n} - n|/n]$ asymptotically equivalent to an oracle constructed over $S = \bigcup_{i=1}^k S_i$ in a single pass. The operation $\cup$ must be commutative and associative. For fault tolerance under at-least-once delivery semantics, the operation must additionally be idempotent.
\end{definition}

Commutativity and associativity guarantee that the final result of combining a collection of oracles is independent of the order of aggregation. A system can perform reductions in parallel, in a hierarchical tree structure, or in any arbitrary sequence without requiring coordination or ordering constraints.

Idempotency ensures that re-combining a duplicate oracle does not alter the aggregate state. This simplifies fault-tolerance protocols. A distributed aggregation system can employ message delivery with at-least-once semantics, obviating the need for more complex exactly-once delivery mechanisms. If a network failure causes a partial estimate to be retransmitted and merged multiple times, the idempotency property guarantees that the final result remains correct.

Collectively, these properties enable the non-additive \texttt{COUNT(DISTINCT)} problem to be reframed as a decomposable aggregation problem, analogous to a \texttt{SUM} operation. The communication cost of distributed cardinality estimation becomes proportional to the size of the compact oracle representation, not to the dataset cardinality.

\subsection{HyperLogLog} \label{subsec:hyperloglog}

The HyperLogLog sketch satisfies the conditions  established in \ref{def:composability}. Its union operation executes as a register-wise maximum, requiring time proportional only to the number of registers $m$, independent of dataset cardinality. This operation is commutative, associative, and idempotent. The property derives from the design of the sketch as a fixed-size summary that encodes cardinality information through the maximum observable in each of $m$ independent stochastic experiments.

\subsubsection{Stochastic Averaging and the Harmonic Mean}

We sum up the core mechanics of the HyperLogLog algorithm as introduced by Flajolet et al. \cite{Flajolet2007HyperLogLogTA}. The algorithm maps elements to an array of registers and records an observable based on the bit patterns of their hash values.

\begin{definition}[HyperLogLog Sketch] \label{def:hll-sketch}
    A HyperLogLog sketch consists of an array $M$ of $m=2^p$ registers, where $p$ is a precision parameter. The sketch is associated with a hash function $h: \mathcal{U} \to \{0,1\}^L$ that distributes elements uniformly over the output space and maps elements from a universe $\mathcal{U}$ to binary strings of sufficient length $L$. Each register $M[j]$ is initialized to 0.
\end{definition}

The algorithm processes a multiset by updating this array of registers. For each element, a single hash value is used to select a register and compute an update value. This process, known as stochastic averaging, uses one hash function to simulate $m$ independent experiments, with each register representing one experiment.

To insert an element $v \in S$, we first compute its hash value $x = h(v)$. The first $p$ bits of $x$ are used to determine a register index $j \in \{0, \dots, m-1\}$. Let $w$ be the remaining $L-p$ bits of $x$. We then compute an observable $\rho(w) = 1 + \text{lz}(w)$, where $\text{lz}(w)$ denotes the number of leading zeros in $w$. The corresponding register is updated according to the rule $M[j] := \max(M[j], \rho(w))$. After all elements have been processed, each register $M[j]$ contains the maximum $\rho$ value observed for any element mapped to that register.

The final cardinality estimate is derived from the state of these registers. The algorithm uses the harmonic mean of the values $2^{M[j]}$, which is less sensitive to outlier values than the arithmetic or geometric mean.

\begin{proposition}[HyperLogLog Estimator \cite{Flajolet2007HyperLogLogTA}]
    Given a HyperLogLog sketch with $m$ registers $M$, the cardinality estimate $E$ is computed as:
    \begin{equation*}
        E = \alpha_m m^2 \left( \sum_{j=0}^{m-1} 2^{-M[j]} \right)^{-1}
    \end{equation*}
    where $\alpha_m$ is a bias correction constant that depends on $m$.
\end{proposition}

This mechanism produces a fixed-size sketch of $\Theta(m \log\log n)$ bits that summarizes the cardinality of a set containing $n$ elements. The theoretical relative error of this estimate is approximately $1.04/\sqrt{m}$ \cite{Flajolet2007HyperLogLogTA}.

\subsubsection{Union and Algebraic Properties} \label{subsubsec:hll-union}

The utility of the HyperLogLog sketch in distributed systems derives from its support for an efficient union operation. This operation allows sketches computed on disparate data partitions to be combined to produce a sketch representing the union of the underlying sets.

\begin{definition}[HyperLogLog Union \cite{Heule2013HyperLogLogIP}]
    Let $M_1$ and $M_2$ be two HyperLogLog sketches of $m$ registers, representing sets $S_1$ and $S_2$ respectively. The union sketch, $M_{1 \cup 2}$, representing the set $S_1 \cup S_2$, is computed as the register-wise maximum of the two input sketches:
    \begin{equation*}
        M_{1 \cup 2}[j] := \max(M_1[j], M_2[j]) \quad \forall j \in \{0, \ldots, m-1\}
    \end{equation*}
\end{definition}

\begin{theorem}[Union Equivalence \cite{Heule2013HyperLogLogIP}]
    The sketch $M_{1 \cup 2}$ produced by the union operation is equivalent to the sketch that would have been generated by processing all elements from both $S_1$ and $S_2$ in a single pass.
\end{theorem}

\begin{proof}
    For any element $x \in S_1 \cup S_2$, its contribution is to set $M[j] := \max(M[j], \rho(w))$ where $j = h(x) \bmod m$ and $\rho(w)$ is the observable derived from $x$. In the single-pass case, this update occurs when $x$ is processed. In the merge case, if $x \in S_1$, then $M_1[j] \geq \rho(w)$ after $M_1$ is constructed. When merging, $M_{1 \cup 2}[j] = \max(M_1[j], M_2[j]) \geq \rho(w)$. Since the update rule is $M[j] := \max(M[j], \rho(w))$, the final value $M_{1 \cup 2}[j]$ is the maximum $\rho$ observed for any element in $S_1 \cup S_2$ that maps to register $j$, which is identical to the value obtained in a single pass over $S_1 \cup S_2$.
\end{proof}

\begin{theorem}[Properties of the Union Operation \cite{Heule2013HyperLogLogIP}] \label{thm:hll-union-properties}
    The HyperLogLog union operation is commutative, associative, and idempotent. For any sketches $M_A$, $M_B$, and $M_C$ we have commutativity
    \begin{equation*}
        M_A \cup M_B = M_B \cup M_A
    \end{equation*}
    associativity
    \begin{equation*}
        (M_A \cup M_B) \cup M_C = M_A \cup (M_B \cup M_C)
    \end{equation*}
    and idempotency
    \begin{equation*}
        M_A \cup M_A = M_A
    \end{equation*}


    % \begin{itemize}
    %     \item Commutativity: $M_A \cup M_B = M_B \cup M_A$
    %     \item Associativity: $(M_A \cup M_B) \cup M_C = M_A \cup (M_B \cup M_C)$
    %     \item Idempotency: $M_A \cup M_A = M_A$
    % \end{itemize}
\end{theorem}

The register-wise maximum operation constitutes the binary operation $\cup$ required by \ref{def:composability}. The commutativity and associativity guarantee order-independent aggregation. The idempotency ensures that the operation satisfies the requirements for fault-tolerant distributed reduction.

Commutativity and associativity guarantee that the final result of merging a collection of sketches is independent of the order of aggregation. This allows a system to perform reductions in parallel, in a hierarchical tree structure, or in any arbitrary sequence without requiring coordination or ordering constraints. Idempotency ensures that re-merging a duplicate sketch does not alter the aggregate state. This simplifies fault-tolerance protocols, as it allows for message delivery with at-least-once semantics, obviating the need for more complex exactly-once delivery mechanisms for sketch aggregation \cite{Heule2013HyperLogLogIP}.

Collectively, these properties enable the non-additive \texttt{COUNT(DISTINCT)} problem to be reframed as a decomposable aggregation problem, analogous to a \texttt{SUM}. This allows the computation to be parallelized with minimal communication overhead \cite{Zaharia2012ResilientDD}.

\subsection{Architectural Implications}

These properties, however, describe the logical behavior of sketch composition, not the physical realities of executing such compositions over millions of instances in a resource-constrained environment. When applied to common analytical workloads, such as large-scale group-by aggregations, the standard HLL algorithm encounters limitations related to estimation accuracy at low cardinalities and communication overhead not apparent from its initial formulation.

\subsubsection{Challenges at Scale}

We analyze two primary challenges that arise when deploying the standard HyperLogLog algorithm in a data-parallel execution environment. Both are consequences of the data distribution patterns common in large-scale analytical queries.

First, analytical workloads frequently involve a high degree of cardinality skew. In a distributed \texttt{GROUP BY} operation executed over billions of records, the resulting groups often follow a power-law distribution. A small number of groups may have very high cardinalities, but the vast majority of groups will contain very few unique elements \cite{Heule2013HyperLogLogIP}. The original HLL estimator exhibits significant systematic bias and a larger relative error for these low-cardinality sets, which compromises the accuracy of the estimates for the long tail of the distribution.

Second, the aggregation of a massive number of sketches introduces a new form of communication overhead. In a distributed \texttt{GROUP BY} query, the map phase generates a separate HLL sketch for each group on each worker node. While each individual sketch is small, the total number of sketches can be in the millions or billions. The subsequent shuffle phase must serialize, transmit, and deserialize this large volume of small objects. This process can itself become a network and CPU bottleneck, undermining the efficiency gains achieved by avoiding the shuffle of raw data \cite{Heule2013HyperLogLogIP}.
%

\subsubsection{HyperLogLog++}

The HyperLogLog++ algorithm introduces a series of engineering refinements to the standard HLL to address the challenges of accuracy and communication overhead at scale \cite{Heule2013HyperLogLogIP}.

First, to address the systematic bias at low cardinalities, it replaces the small-range correction with a bias correction mechanism derived from empirical data. The system uses a pre-computed table of bias values obtained from simulations to correct the raw estimate, improving accuracy for the low-cardinality groups prevalent in skewed analytical workloads \cite{Heule2013HyperLogLogIP}.

Second, it forces the use of a 64-bit hash function. This change expands the hash space to make collisions statistically insignificant even for cardinalities exceeding billions, thereby eliminating the need for the large-range correction formula of the original algorithm and simplifying the implementation \cite{Heule2013HyperLogLogIP}.

The most significant architectural refinement is the introduction of a dual-representation sketch that adapts its storage format based on cardinality. This design directly addresses the communication overhead associated with shuffling a massive number of sketches.

\begin{definition}[Sparse and Dense Representations \cite{Heule2013HyperLogLogIP}]
    A HyperLogLog++ sketch can exist in one of two formats. A \texttt{dense} format, which is the standard HLL representation: a fixed-size array of $m$ registers. A \texttt{sparse} format, used for low cardinalities, which stores only the non-zero registers as a variable-length encoded, sorted list of (index, value) pairs. The sketch begins in the \texttt{sparse} format and transitions to the \texttt{dense} format when the encoded byte size of the sparse representation exceeds the byte size of the dense format.
\end{definition}

The sparse format reduces the memory footprint of these numerous low-cardinality sketches by an order of magnitude or more. This reduction in size directly translates to a commensurate reduction in the volume of data that must be serialized, transmitted during the shuffle phase, and deserialized by reducer nodes \cite{Heule2013HyperLogLogIP}.

\subsection{Architectures and Implementations}

The properties of HyperLogLog sketches, particularly their efficient and composable union operation, enable their integration into diverse distributed computing architectures. The choice of architecture is typically dictated by the nature of the data workload, whether batch or streaming, and the desired query patterns.

\subsubsection{Architectural Patterns} \label{subsubsec:hll-patterns}

We analyze three canonical patterns for distributed HLL computation. The first is the MapReduce paradigm. In this model, an input dataset is partitioned across worker nodes. The map stage generates local HLL sketches for each data partition, creating a separate sketch for each grouping key in a group-by operation. The system then shuffles these compact sketches to reducer nodes. In the reduce stage, each reducer applies the union operation to merge the sketches it receives for a given key. The correctness of this stage, where sketches may arrive in an arbitrary order, is guaranteed by the commutativity and associativity of the union operation \cite{Zaharia2012ResilientDD}.

The second pattern is designed for real-time aggregation in stream processing systems. This architecture partitions an unbounded data stream into logical windows, such as fixed time intervals or element counts. For each active window, an operator maintains a local HLL sketch. The composability of sketches allows for efficient computation over combined windows. To compute the cardinality over a time interval $T$ composed of non-overlapping sub-intervals $t_1, \ldots, t_k$, the system retrieves the corresponding sketches $M_{t_1}, \ldots, M_{t_k}$ and computes the final sketch as $M_T = \bigcup_{i=1}^{k} M_{t_i}$. For overlapping windows, the system must maintain separate sketches for each window to avoid double-counting elements. The cost of this operation is independent of the raw data volume within the interval, enabling efficient, stateful stream processing \cite{Huang2016TowardHD}.

The third pattern is a hierarchical, or tree-based, aggregation model, well-suited for systems with a natural data hierarchy \cite{Wilms2021TreebasedNA}. In this model, leaf nodes generate sketches from raw data streams. Intermediate nodes in the hierarchy periodically collect and merge the sketches from their children to create higher-level aggregates. The validity of this hierarchical model rests upon the associativity of the HLL union, which ensures that the final global sketch is identical regardless of the tree's depth or the order of intermediate merge operations. This pattern enables multi-resolution querying, where queries for cardinality at different levels of the hierarchy can be answered by accessing pre-aggregated sketches without reprocessing data from lower levels.

\begin{remark}
    In hierarchical aggregation, the relative error of the global estimate remains $O(1/\sqrt{m})$ independent of tree depth, as each merge produces a sketch equivalent to single-pass construction over the union of the underlying sets.
\end{remark}

% These three patterns instantiate the hierarchical aggregation architecture analyzed in the framework (\S~\ref{subsubsec:composability-hierarchy}). The framework establishes that composability-optimized structures enable parallel reduction through tree-based merge operations. The MapReduce pattern implements a two-level hierarchy (workers and reducers). The stream processing pattern implements temporal hierarchies (window aggregation). The tree-based pattern implements spatial hierarchies (organizational or geographical aggregation levels). All three patterns rely on the properties established in Theorem~\ref{thm:hll-union-properties} to guarantee correctness under arbitrary merge orderings.

\subsubsection{Implementation Case Studies}

The patterns enabled by HyperLogLog are realized in a spectrum of production systems, ranging from low-level data structures requiring explicit user orchestration to high-level, transparent query optimizations managed entirely by a database engine. We analyze three representative implementations.

The Redis implementation \cite{redis-hll} provides HLL as a primitive data type. It exposes three core commands: \texttt{PFADD} to insert elements, \texttt{PFCOUNT} to retrieve the cardinality estimate, and \texttt{PFMERGE} to perform the union of multiple sketches. An HLL sketch is represented as a fixed-size string. This facilitates a client-driven model for distributed computation: an application can retrieve a serialized sketch from one node, transmit it over the network, and merge it with another sketch on a different node by invoking \texttt{PFMERGE}. In this model, the orchestration of the distributed aggregation is the explicit responsibility of the application logic.

In contrast, distributed SQL query engines such as Presto and Google BigQuery integrate HLL as a high-level, transparent optimization within the query execution layer. These systems expose the functionality through an aggregate function, typically \texttt{APPROX\_DISTINCT}. When executing a distributed query, the query planner automatically pushes the sketch creation stage to the worker nodes. The workers compute HLL++ sketches, often utilizing both sparse and dense representations to optimize for memory and network usage. In the final stage, these partial sketches are shuffled to a reducer or coordinator node, which performs a final merge aggregation. The entire process is managed by the query optimizer, abstracting the underlying mechanics from the user \cite{Heule2013HyperLogLogIP}. BigQuery further provides a multi-stage API (\texttt{HLL\_COUNT.INIT}, \texttt{HLL\_COUNT.MERGE}) that allows intermediate, partially aggregated sketches to be materialized in tables, enabling efficient incremental rollups \cite{Tao2021Expected1O}.

The highest level of abstraction is represented by systems like the Citus extension for PostgreSQL. This system provides transparent distributed cardinality estimation through automatic query rewriting. An administrator can enable approximate counting via a configuration parameter. Subsequently, the distributed query planner automatically rewrites standard \texttt{COUNT(DISTINCT)} queries to use an HLL-based implementation. It pushes sketch aggregation down to the worker nodes and merges the intermediate sketches on the coordinator node to produce the final result. From the user's perspective, the use of HLL is a transparent physical-layer optimization managed entirely by the database \cite{Harmouch2017CardinalityEA}.

\subsection{Challenges in Production Environments}

The deployment of probabilistic data structures in production systems introduces several difficulties. We will analyze three such challenges: ensuring fault tolerance efficiently, providing security against adversarial manipulation, and managing privacy in federated contexts.

\subsubsection{Fault Tolerance and Approximate Recovery} \label{subsubsec:hll-fault-tolerance}

In data-parallel frameworks, fault tolerance for stateless transformations is achieved by recomputing lost partitions from their lineage graph \cite{Zaharia2012ResilientDD}. For stateful operations, such as maintaining HLL sketches over streaming data, this approach is insufficient. Traditional fault tolerance for stateful operators requires periodic checkpointing of the entire state, which incurs significant I/O overhead and compromises performance \cite{Huang2016TowardHD}.

An alternative approach, termed approximate fault tolerance, leverages the inherently probabilistic nature of the data structure to create more efficient recovery mechanisms. Instead of checkpointing at fixed intervals, this model triggers a backup only when the potential error introduced by a failure would exceed a user-defined threshold. This is managed by monitoring the divergence between the current in-memory state and its last persisted backup \cite{Huang2016TowardHD}.

\begin{definition}[State Divergence \cite{Huang2016TowardHD}]
    Let $M_{cur}$ be the current state of an HLL sketch and $M_{back}$ be its most recent persisted state. The state divergence $D(M_{cur}, M_{back})$ quantifies the difference between the two sketches. For HLL, this can be defined as:
    \begin{equation*}
        D(M_{cur}, M_{back}) = |\text{Estimate}(M_{cur}) - \text{Estimate}(M_{back})|
    \end{equation*}
    or alternatively as the sum of register-wise differences:
    \begin{equation*}
        D(M_{cur}, M_{back}) = \sum_{j=0}^{m-1} |M_{cur}[j] - M_{back}[j]|
    \end{equation*}
\end{definition}

It executes a backup operation only when the state divergence $D$ exceeds a pre-configured threshold $\Theta$, or when the number of unprocessed items since the last backup exceeds another threshold $L$ \cite{Huang2016TowardHD}.

\begin{proposition}[Bounded Error Under Failures \cite{Huang2016TowardHD}]
    In the AF-Stream model, let a worker experience $k$ failures. The total accumulated divergence $D_{total}$ between the actual state and the ideal, failure-free state satisfies $D_{total} \leq k \cdot \max(\Theta, c \cdot L)$ where $\Theta$ and $L$ are the user-specified divergence and item count thresholds, and $c$ is an algorithm-specific constant. The bound grows linearly with $k$ but is controlled by the threshold parameters.
\end{proposition}

% This model demonstrates a principled trade-off, where a system accepts a small, provably bounded increase in estimation error in the rare event of a failure in exchange for a significant reduction in the constant overhead of checkpointing.

\subsubsection{Security Against Adversarial Manipulation} \label{subsubsec:hll-security}

The probabilistic nature of HyperLogLog, while providing efficiency, also introduces security vulnerabilities. We analyze a specific class of adversarial action, the cardinality inflation attack, where a malicious actor seeks to deliberately produce an arbitrarily large cardinality estimate from a small set of crafted inputs \cite{Reviriego2020HyperLogLogS}.

The attack leverages the internal structure of the HLL algorithm. An attacker can empirically discover input elements that produce high $\rho$ values for each of the $m$ registers. Even without knowledge of the specific hash function, an attacker can construct a malicious set by repeatedly generating random inputs and observing their effect on a local HLL instance. By collecting just $m$ such inputs, one targeting each register, an attacker can construct a set that populates every register with a high value. When this small set is inserted into the target HLL sketch, the harmonic mean calculation yields a massive, artificially inflated cardinality estimate. This vulnerability has been demonstrated to be effective against production implementations in systems such as Presto and Redis \cite{Reviriego2020HyperLogLogS}.

We consider two classes of mitigation strategies. The first involves cryptographic techniques, such as applying a secret, system-specific salt to the hash function. This approach makes it computationally difficult for an external attacker to craft malicious inputs offline. However, this strategy breaks the mergeability property for sketches that use different salts, limiting its applicability in systems that rely on sketch union for aggregation across different administrative domains \cite{Reviriego2020HyperLogLogS}.

The second class of strategies involves an aggregation mechanisms designed to operate in the presence of Byzantine adversaries. Instead of a simple \texttt{max} operation during a merge, or a simple harmonic mean for estimation, the system can employ aggregation rules that are resilient to outliers. Techniques from the field of Byzantine-resilient distributed optimization, such as the Geometric Median or aggregation rules like Krum, are designed to filter out malicious updates \cite{Le2025ByzantineRF}. The Geometric Median, for instance, finds a point that minimizes the sum of Euclidean distances to a set of input vectors, making it inherently robust to arbitrarily corrupted inputs. Applying such aggregators to the vector of HLL registers during a merge operation is an active area of research for building secure, Byzantine-fault-tolerant analytical systems \cite{Kuwaranancharoen2023OnTG}.

\subsubsection{Privacy Implications in Federated Systems} \label{subsubsec:hll-privacy}

While HLL sketches provide a compact summary that does not store raw data elements, their use in federated environments introduces privacy considerations. The act of sharing a sketch, even if it contains only anonymized hash-derived values, can still leak information about the underlying dataset, particularly concerning unique or rare elements \cite{Le2025ByzantineRF}.

In a federated query system, an adversary who compromises the central aggregator or observes the transmitted sketches could perform linkage attacks. If a hospital finds that only one patient satisfies the criteria for a query, sharing an HLL sketch of this single-element set still reveals information. An attacker could potentially infer properties of this unique patient by correlating the sketch with a known background population \cite{Le2025ByzantineRF}. To quantify this risk, we can apply the concept of k-anonymity to the buckets of an HLL sketch.

\begin{definition}[Bucket k-Anonymity \cite{Le2025ByzantineRF}]
    Let $B$ be a set of elements matching a query, drawn from a background population $A$. Consider an HLL bucket $j$ containing value $M[j] = v$. This bucket is k-anonymous if $|\{x \in A : h(x) \bmod m = j \text{ and } \rho(x) = v\}| \geq k$, meaning at least $k$ distinct elements from $A$ could have produced this bucket value. An attacker observing the sketch cannot determine with certainty which of these $k$ elements contributed to the bucket's value.
\end{definition}

The level of k-anonymity provided by an HLL sketch is a function of the sketch parameters, the size of the background population, and the size of the query set. Research has shown that it is possible to compute theoretical bounds on the expected number of buckets in a sketch that are not k-anonymous. HLL sketches are particularly effective for rare diseases (where the ratio of query set size to background population size is low). This is in marked contrast to sharing exact counts, where a count of one immediately breaks anonymity \cite{Le2025ByzantineRF}. This allows us to manage the trade-off between the accuracy of the cardinality estimate, which increases with the number of buckets $m$, and the privacy risk, as a larger $m$ may decrease the expected k-anonymity of each bucket.

\subsection{Limitations and Variants}


The value of a probabilistic data structure is determined not only by its space-error characteristics but also by the set of operations it supports. The HyperLogLog algorithm is highly specialized for cardinality estimation and set union. While its merge operation is fundamental for distributed aggregation, it does not natively support other operations such as intersection or set difference.

Alternative data structures, such as MinHash and its derivatives like Theta Sketches, provide a richer set of supported operations. These sketches are designed to estimate the Jaccard similarity between sets, from which the sizes of their union, intersection, and difference can be derived. However, this increased functionality comes at a cost. For the primary task of cardinality estimation, HLL is substantially more space-efficient than MinHash-based sketches providing the same estimation accuracy \cite{Ertl2021SetSketchFT}.

% This distinction represents a trade-off between functional generality and architectural efficiency. HLL is architecturally optimized for the specific problem of distributed, union-based cardinality aggregation. Its design provides a near-optimal solution for this problem class. For analytical queries that require measuring the overlap between sets, a different class of data structures is necessary. The choice between these algorithms is therefore not an optimization detail but a primary architectural decision dictated by the functional requirements of the system.

\subsubsection{Variants for Cardinality Estimation}

% Research in cardinality estimation continues to refine the space-error trade-off and explore new architectural possibilities. Recent algorithmic developments demonstrate that further space efficiencies are achievable while preserving the essential algebraic properties required for distributed deployment.

The UltraLogLog algorithm modifies the underlying register structure to encode more information per bit \cite{Ertl2023UltraLogLogAP}. It generalizes the HyperLogLog register by partitioning its bits. A set of $q$ bits stores the maximum update value $u_i$ seen so far, while a set of $d$ additional bits acts as a compact summary, indicating whether updates with values $u_i-1, \dots, u_i-d$ have occurred. This design, particularly for the ULL configuration with $q=6$ and $d=2$, fits into a single 8-bit register. It achieves a theoretical memory-variance product approximately 28\% lower than a standard 6-bit HLL sketch while remaining commutative, idempotent, and mergeable.

The HyperLogLogLog algorithm\footnote{Three log scientist \url{https://curiouscoding.nl/posts/three-log-scientist/}} achieves compression through a sparse-dense architecture based on value offsets \cite{Karppa2022HyperLogLogLogCE}.

\begin{definition}[HyperLogLogLog Sketch Structure \cite{Karppa2022HyperLogLogLogCE}]
    A HyperLogLogLog sketch is a 3-tuple $(S, M, B)$ where $B$ is a base value, representing a lower bound for the majority of register values.

    $M$ is a dense array storing small-width integer offsets for registers whose values are close to $B$. For a register $j$ where $B \leq M'[j] < B + 2^k$, its offset $M[j] = M'[j] - B$ is stored using $k$ bits.

    $S$ is a sparse associative array storing pairs $(j, M'[j])$ for outlier registers whose values fall outside the dense offset range. $S$ is a sparse associative array storing pairs $(j, M'[j])$ for outlier registers whose values fall outside the dense offset range.

    The base value $B$ is chosen dynamically to minimize the total representation size $k \cdot m + |S| \cdot (\log m + \log w)$; see \cite{Karppa2022HyperLogLogLogCE} for the selection algorithm.
\end{definition}

This structure adapts to the statistical distribution of register values, using compact offsets for the common case and a sparse map for exceptions. The resulting sketch requires approximately $m \log_2\log_2\log_2 m$ bits for large $n$, albeit with a potential increase in the computational complexity of updates \cite{Karppa2022HyperLogLogLogCE}.

Several open problems remain at the intersection of cardinality estimation and distributed systems. The development of robust aggregation protocols that are resilient to Byzantine adversaries is an active area of research. Securing sketch-based systems requires moving beyond simple merge logic to robust aggregation rules that can tolerate corrupted inputs. Furthermore, integrating formal privacy guarantees with cardinality estimation sketches is necessary for enabling collaborative analytics on sensitive data. Finally, new architectural paradigms such as sketch disaggregation, where a single logical sketch is fragmented and distributed across network hardware, present a frontier for research into ultra-low-latency, in-network analytics \cite{Langlet2025SketchDA}.

% ==============================================================================
% FINE DELLA SEZIONE
% ==============================================================================
